\iffalse
!!!!!!!!!!!!!!!!!!!!!!!!!!!!!!!!!!!!!
! Do not modify this file directly  !
! This file has been auto-generated !
!!!!!!!!!!!!!!!!!!!!!!!!!!!!!!!!!!!!!
\fi

%% ODER: format ==         = "\mathrel{==}"
%% ODER: format /=         = "\neq "
%
%
\makeatletter
\@ifundefined{lhs2tex.lhs2tex.sty.read}%
  {\@namedef{lhs2tex.lhs2tex.sty.read}{}%
   \newcommand\SkipToFmtEnd{}%
   \newcommand\EndFmtInput{}%
   \long\def\SkipToFmtEnd#1\EndFmtInput{}%
  }\SkipToFmtEnd

\newcommand\ReadOnlyOnce[1]{\@ifundefined{#1}{\@namedef{#1}{}}\SkipToFmtEnd}
\usepackage{newtxmath}
\usepackage{stmaryrd}
\DeclareFontFamily{OT1}{cmtex}{}
\DeclareFontShape{OT1}{cmtex}{m}{n}
  {<5><6><7><8>cmtex8
   <9>cmtex9
   <10><10.95><12><14.4><17.28><20.74><24.88>cmtex10}{}
\DeclareFontShape{OT1}{cmtex}{m}{it}
  {<-> ssub * cmtt/m/it}{}
\newcommand{\texfamily}{\fontfamily{cmtex}\selectfont}
\DeclareFontShape{OT1}{cmtt}{bx}{n}
  {<5><6><7><8>cmtt8
   <9>cmbtt9
   <10><10.95><12><14.4><17.28><20.74><24.88>cmbtt10}{}
\DeclareFontShape{OT1}{cmtex}{bx}{n}
  {<-> ssub * cmtt/bx/n}{}
\newcommand{\tex}[1]{\text{\texfamily#1}}	% NEU

\newcommand{\Sp}{\hskip.33334em\relax}


\newcommand{\Conid}[1]{\mathit{#1}}
\newcommand{\Varid}[1]{\mathit{#1}}
\newcommand{\anonymous}{\kern0.06em \vbox{\hrule\@width.5em}}
\newcommand{\plus}{\mathbin{+\!\!\!+}}
\newcommand{\bind}{\mathbin{>\!\!\!>\mkern-6.7mu=}}
\newcommand{\rbind}{\mathbin{=\mkern-6.7mu<\!\!\!<}}% suggested by Neil Mitchell
\newcommand{\sequ}{\mathbin{>\!\!\!>}}
\renewcommand{\leq}{\leqslant}
\renewcommand{\geq}{\geqslant}
\usepackage{polytable}

%mathindent has to be defined
\@ifundefined{mathindent}%
  {\newdimen\mathindent\mathindent\leftmargini}%
  {}%

\def\resethooks{%
  \global\let\SaveRestoreHook\empty
  \global\let\ColumnHook\empty}
\newcommand*{\savecolumns}[1][default]%
  {\g@addto@macro\SaveRestoreHook{\savecolumns[#1]}}
\newcommand*{\restorecolumns}[1][default]%
  {\g@addto@macro\SaveRestoreHook{\restorecolumns[#1]}}
\newcommand*{\aligncolumn}[2]%
  {\g@addto@macro\ColumnHook{\column{#1}{#2}}}

\resethooks

\newcommand{\onelinecommentchars}{\quad-{}- }
\newcommand{\commentbeginchars}{\enskip\{-}
\newcommand{\commentendchars}{-\}\enskip}

\newcommand{\visiblecomments}{%
  \let\onelinecomment=\onelinecommentchars
  \let\commentbegin=\commentbeginchars
  \let\commentend=\commentendchars}

\newcommand{\invisiblecomments}{%
  \let\onelinecomment=\empty
  \let\commentbegin=\empty
  \let\commentend=\empty}

\visiblecomments

\newlength{\blanklineskip}
\setlength{\blanklineskip}{0.66084ex}

\newcommand{\hsindent}[1]{\quad}% default is fixed indentation
\let\hspre\empty
\let\hspost\empty
\newcommand{\NB}{\textbf{NB}}
\newcommand{\Todo}[1]{$\langle$\textbf{To do:}~#1$\rangle$}

\EndFmtInput
\makeatother
%
%
%
%
%
%
% This package provides two environments suitable to take the place
% of hscode, called "plainhscode" and "arrayhscode". 
%
% The plain environment surrounds each code block by vertical space,
% and it uses \abovedisplayskip and \belowdisplayskip to get spacing
% similar to formulas. Note that if these dimensions are changed,
% the spacing around displayed math formulas changes as well.
% All code is indented using \leftskip.
%
% Changed 19.08.2004 to reflect changes in colorcode. Should work with
% CodeGroup.sty.
%
\ReadOnlyOnce{polycode.fmt}%
\makeatletter

\newcommand{\hsnewpar}[1]%
  {{\parskip=0pt\parindent=0pt\par\vskip #1\noindent}}

% can be used, for instance, to redefine the code size, by setting the
% command to \small or something alike
\newcommand{\hscodestyle}{}

% The command \sethscode can be used to switch the code formatting
% behaviour by mapping the hscode environment in the subst directive
% to a new LaTeX environment.

\newcommand{\sethscode}[1]%
  {\expandafter\let\expandafter\hscode\csname #1\endcsname
   \expandafter\let\expandafter\endhscode\csname end#1\endcsname}

% "compatibility" mode restores the non-polycode.fmt layout.

\newenvironment{compathscode}%
  {\par\noindent
   \advance\leftskip\mathindent
   \hscodestyle
   \let\\=\@normalcr
   \let\hspre\(\let\hspost\)%
   \pboxed}%
  {\endpboxed\)%
   \par\noindent
   \ignorespacesafterend}

\newcommand{\compaths}{\sethscode{compathscode}}

% "plain" mode is the proposed default.
% It should now work with \centering.
% This required some changes. The old version
% is still available for reference as oldplainhscode.

\newenvironment{plainhscode}%
  {\hsnewpar\abovedisplayskip
   \advance\leftskip\mathindent
   \hscodestyle
   \let\hspre\(\let\hspost\)%
   \pboxed}%
  {\endpboxed%
   \hsnewpar\belowdisplayskip
   \ignorespacesafterend}

\newenvironment{oldplainhscode}%
  {\hsnewpar\abovedisplayskip
   \advance\leftskip\mathindent
   \hscodestyle
   \let\\=\@normalcr
   \(\pboxed}%
  {\endpboxed\)%
   \hsnewpar\belowdisplayskip
   \ignorespacesafterend}

% Here, we make plainhscode the default environment.

\newcommand{\plainhs}{\sethscode{plainhscode}}
\newcommand{\oldplainhs}{\sethscode{oldplainhscode}}
\plainhs

% The arrayhscode is like plain, but makes use of polytable's
% parray environment which disallows page breaks in code blocks.

\newenvironment{arrayhscode}%
  {\hsnewpar\abovedisplayskip
   \advance\leftskip\mathindent
   \hscodestyle
   \let\\=\@normalcr
   \(\parray}%
  {\endparray\)%
   \hsnewpar\belowdisplayskip
   \ignorespacesafterend}

\newcommand{\arrayhs}{\sethscode{arrayhscode}}

% The mathhscode environment also makes use of polytable's parray 
% environment. It is supposed to be used only inside math mode 
% (I used it to typeset the type rules in my thesis).

\newenvironment{mathhscode}%
  {\parray}{\endparray}

\newcommand{\mathhs}{\sethscode{mathhscode}}

% texths is similar to mathhs, but works in text mode.

\newenvironment{texthscode}%
  {\(\parray}{\endparray\)}

\newcommand{\texths}{\sethscode{texthscode}}

% The framed environment places code in a framed box.

\def\codeframewidth{\arrayrulewidth}
\RequirePackage{calc}

\newenvironment{framedhscode}%
  {\parskip=\abovedisplayskip\par\noindent
   \hscodestyle
   \arrayrulewidth=\codeframewidth
   \tabular{@{}|p{\linewidth-2\arraycolsep-2\arrayrulewidth-2pt}|@{}}%
   \hline\framedhslinecorrect\\{-1.5ex}%
   \let\endoflinesave=\\
   \let\\=\@normalcr
   \(\pboxed}%
  {\endpboxed\)%
   \framedhslinecorrect\endoflinesave{.5ex}\hline
   \endtabular
   \parskip=\belowdisplayskip\par\noindent
   \ignorespacesafterend}

\newcommand{\framedhslinecorrect}[2]%
  {#1[#2]}

\newcommand{\framedhs}{\sethscode{framedhscode}}

% The inlinehscode environment is an experimental environment
% that can be used to typeset displayed code inline.

\newenvironment{inlinehscode}%
  {\(\def\column##1##2{}%
   \let\>\undefined\let\<\undefined\let\\\undefined
   \newcommand\>[1][]{}\newcommand\<[1][]{}\newcommand\\[1][]{}%
   \def\fromto##1##2##3{##3}%
   \def\nextline{}}{\) }%

\newcommand{\inlinehs}{\sethscode{inlinehscode}}

% The joincode environment is a separate environment that
% can be used to surround and thereby connect multiple code
% blocks.

\newenvironment{joincode}%
  {\let\orighscode=\hscode
   \let\origendhscode=\endhscode
   \def\endhscode{\def\hscode{\endgroup\def\@currenvir{hscode}\\}\begingroup}
   %\let\SaveRestoreHook=\empty
   %\let\ColumnHook=\empty
   %\let\resethooks=\empty
   \orighscode\def\hscode{\endgroup\def\@currenvir{hscode}}}%
  {\origendhscode
   \global\let\hscode=\orighscode
   \global\let\endhscode=\origendhscode}%

\makeatother
\EndFmtInput
%
%


%
%
%
%
%
\ReadOnlyOnce{lineno.fmt}%

\usepackage{etoolbox}

\makeatletter

\newtoggle{lineno@column}
\newtoggle{lineno@print}
\newcounter{lineno@num}
\newlength\linenowidth

\newcommand\linenoFormat{\footnotesize}

\newcommand\linenoSetup{%
  \iftoggle{lineno@column}{%
    \column{lineno}{@{}r}%
    \column{B}{}
  }{}%
}

\newcommand\linenoStart{%
  \refstepcounter{lineno@num}%
  \iftoggle{lineno@print}{%
    \fromto{lineno}{B}{\makebox[\linenowidth][r]{\linenoFormat\arabic{lineno@num}}}%
  }{}%
}

\newcommand\linenoon[1][]{%
  #1\toggletrue{lineno@column}%
  #1\toggletrue{lineno@print}%
}

\newcommand\linenooff[1][]{%
  #1\togglefalse{lineno@column}%
  #1\togglefalse{lineno@print}%
}

\newcommand\linenoskip[1][]{%
  #1\toggletrue{lineno@column}%
  #1\togglefalse{lineno@print}%
}

\newcommand\linenoreset[1][1]{%
  \defcounter{lineno@num}{#1 - 1}%
}

\EndFmtInput












\section{Introduction}
\label{sec:introduction}

Concurrent programs decompose computation into activities that
communicate across processes, cores, or machines. Such communication is
governed by protocols, and these protocols are easy to get wrong: a
message may be missing, a choice may be unexpected, or an action may
occur in the wrong state. A type discipline for concurrency should
therefore specify both the data exchanged by components and the order
in which those exchanges may occur.

Session types
\cite{DBLP:conf/concur/Honda93,HondaVK98,Gay_Vasconcelos_2025}
address precisely this problem: they are a type discipline for
structured communication
protocols. Their core features are typed, bidirectional message
exchange and typed choices. They ensure protocol fidelity as well as type soundness,
for pi calculus as well as for concurrent lambda calculus. There is a significant body of
literature on session types, which is well appreciated in
survey articles and books
\cite{DBLP:journals/csur/HuttelLVCCDMPRT16,
  gay17:_behav_types,
  DBLP:journals/ftpl/AnconaBB0CDGGGH16,
  Gay_Vasconcelos_2025}.
While much recent work concentrates on multi-party session types
\cite{DBLP:journals/jacm/HondaYC16}, this work is solely concerned
with binary session types, which govern two-party protocols.

Session types support infinite and unbounded behaviors through
recursive types almost from the beginning \cite{HondaVK98}. Since
then, most work has concentrated on \emph{regular} recursive session
types, where recursion is restricted to tail positions in the
type. For example, here is the type
\ensuremath{\mu  \alpha .\&\{\mkern1mu\relax \CodeLbl{\Conid{Next}}\mathbin{:}\mathord{?}\Conid{Int}.\alpha ,\CodeLbl{\Conid{Quit}}\mathbin{:}\CodeTy{Close}\mkern1mu\}}
of a server that receives a list of integers. The occurrence of
the recursive \ensuremath{\alpha } is syntactically restricted to tail
position because regular session types are built by composing prefixes
using syntax constructors like \ensuremath{\mathord{?}\Conid{Int}.\anonymous }. Context-free session types
\cite{DBLP:conf/icfp/ThiemannV16,DBLP:journals/iandc/AlmeidaMTV22} (CFSTs)
have lifted this restriction by adopting a symmetric composition
operator \ensuremath{\anonymous \mathop{;}\anonymous } instead. This operator
enables tree- and stack-shaped communication protocols, but it also
makes type conversion non-trivial because session types must be
identified modulo the algebraic laws of sequential composition.

Borrowing has received much attention recently due to its adoption in
the Rust language \cite{team24:_rust_progr_languag}. It refers to mechanisms that
grant temporary access to a resource without transferring ownership of
that resource. Originally conceived for imperative programming,
\citet{DBLP:journals/pacmpl/SaffrichS0V25} present a session type
system with a notion of borrowing.
In this system, a borrow consumes a prefix of a session and leaves a residual suffix, with the
typing discipline internalizing the linear update of the channel.
Their calculus BGV is based on regular
session types \cite{DBLP:conf/esop/LindleyM15} and its semantics is
given by a translation to a calculus from the literature
\cite{DBLP:journals/lmcs/KokkeD23}. However, the approach of the BGV paper has a number of
drawbacks. As the source calculus has no direct semantics, it is
difficult to perform optimizations at the source level. Moreover, the
translation is uniform in the sense that it does not distinguish
between local borrows that remain in the same process and remote
borrows that are passed to other processes. This uniformity leads to
inefficiency because local borrows can elide the synchronization
machinery needed for remote borrows. 

This work promotes the ideas of borrowing for regular session types
\cite{DBLP:journals/pacmpl/SaffrichS0V25} to context-free session
types with borrowing, a calculus we call \thecalculus, both
simplifying and improving on that work. While CFSTs'
symmetric composition enables expressive protocols, it does not
liberate functional session type programming from resource-passing
style. \Cref{sec:borr-from-cont} develops this point using tree
serialization as an example: the plain CFST program
threads the channel explicitly and relies on polymorphic recursion,
whereas the \thecalculus version uses local borrowing to consume one
prefix of the channel after another. \Cref{sec:local-remote-borr} then
uses a streaming HTML renderer to show why remote borrowing is
different: once a borrow crosses a process boundary, evaluation order
alone no longer sequences the borrowed prefix and the residual
session, so synchronization is needed.

The tree example also exposes the core technical issue. The move from regular 
session types to CFSTs makes borrowing technically more delicate. In a
regular session type, the prefix consumed by a borrow is exposed by
the syntactic structure of the type as well as by unfolding
recursion. For CFSTs, however, the desired 
prefix may only become visible after reassociating sequential
composition, eliminating \ensuremath{\CodeTy{Skip}}, distributing continuations
through choices, or unfolding recursion.  Thus splitting a borrow from
a CFST must be defined modulo the equational theory of CFSTs and
computing the residual suffix from the original type and the borrowed
prefix is a significant task.

\subsection*{Contributions}

The core conceptual contributions of this paper are as follows.
\begin{itemize}
\item We extend borrowing from regular session types to context-free
  session types. The resulting calculus, \thecalculus, keeps the
  expressiveness of CFSTs while removing much of the
  resource-passing and continuation-polymorphic programming style that
  plain CFSTs impose.
\item We identify session-type splitting as the central operation
  behind borrowing for CFSTs. Splitting must check a residual suffix
  against a session and a borrowed prefix type modulo the equational
  theory of CFSTs.
\item We distinguish local borrowing from remote borrowing. Local
  borrows remain in one process and can be ordered by evaluation and
  typing alone; remote borrows cross process boundaries and therefore
  require explicit synchronization.
\item We give \thecalculus a direct operational semantics. This
  semantics serves as an abstract specification of borrowing, rather
  than defining the source language only by translation.
\end{itemize}

The remaining contributions are mainly technical.
\begin{itemize}
\item We define a declarative type system with ordered contexts inspired by
  Bunched Implications \cite{DBLP:journals/bsl/OHearnP99}, effects, mobility, and
  boundedness.
\item We present a constraint-based algorithmic typing judgment, together with
  soundness and completeness statements against the declarative system.
\item We define a low-level run-time calculus with explicit channel
  representations and synchronization flags, and translate typed
  \thecalculus processes to this calculus.
\item We state the metatheoretic obligations connecting the layers,
  including type soundness for the direct semantics and a weak simulation
  result for the translation.
\item We describe an implementation path for type checking that
  leverages the implementation of FreeST
  \cite{almeida26:_frees_concur_progr_languag_with,DBLP:journals/corr/abs-1904-01284}.
\end{itemize}

The supplement contains omitted rules and auxiliary material;
cross-references into it are marked by an \textup{S-} prefix, such as
\cref{fig:omitted-predicates} or \cref{sec:appendix-types}.

%%% Local Variables:
%%% mode: latex
%%% TeX-master: "../popl2027.tex"
%%% End:
